\section*{Rangkuman}
\addcontentsline{toc}{section}{Rangkuman}

Pada bab ini telah dibahas:
\begin{itemize}[leftmargin=*]
  \item definisi formal alfabet, string, bahasa, dan tata bahasa;
  \item derivasi, bentuk kalimat, ekivalensi, dan bentuk normal;
  \item hierarki Chomsky Tipe 3 sampai Tipe 0;
  \item otomata hingga, otomata pushdown, otomata terbatas linear,
        dan mesin Turing;
  \item tabel pemetaan Tipe 3--0 dan uji klasifikasi lima grammar;
  \item implikasi hierarki bagi front-end kompilator.
\end{itemize}

Bab V memulai perjalanan implementasi dengan analisis leksikal. Kita
akan melihat bagaimana ekspresi reguler dipakai untuk mendefinisikan
token, dan bagaimana teori Tipe 3 berubah menjadi \eng{lexer} yang
bekerja.
